% =====================================================================
%  ch09a.tex  —— 第 9 章 二维规范理论（原书 p.119–140 / PDF 128–149，前半）
%  译自 植田一石《数物系のためのシンプレクティック幾何学入門》SGC 118
%  本 chunk 含 \chapter{二维规范理论} 与 9.1–9.11；章末回顾见 ch09b。
% =====================================================================

\chapter{二维规范理论}

% [原文 p.119]
如果说环面流形是人造辛流形的代表例，那么天然辛流形的代表例大概就是 Riemann 面上稳定向量丛的模空间了。由 Narasimhan--Seshadri 定理，它与常中心曲率联络的模空间自然地视为同一，构成代数几何、辛几何、规范理论、低维拓扑等相互交织的有趣领域。

% [Tutor Note: Narasimhan–Seshadri 定理（1965）断言紧 Riemann 面上次数与秩互素的稳定全纯向量丛与基本群的中心扩张的不可约酉表示一一对应；Atiyah–Bott 将其重述为无限维辛商，见本章 9.5 节。]

\section{二维的乐园}

居住的空间是 3 维的，时空是 4 维的。这为研究 3 维或 4 维的几何提供了巨大的动力，但为何要研究 2 维的几何呢？

第一个动力是，2 维的模型为 3 维或 4 维的模型提供了玩具。在 2 维中预演，可以成为研究更困难的 3 维、4 维模型的准备。

第二个动力是，存在可近似地看作 2 维的 3 维或 4 维系统。细纤维状的系统的动力学、薄膜上的定常系统，或者在空间上延展但在 1 个方向上均匀的定常系统、在 2 个方向上均匀的系统的动力学等，都属此类。

第三个动力是，模型的数学结构很美，因而要研究它们，至于能否由实验直接检验并不是问题。实际上，只在 2 维才发生的奇迹比比皆是，单从纯粹审美的观点看，研究这些模型也能带来巨大的喜悦。其典型例子是 2 维量子引力理论，它与以 KdV 方程、KP 方程为代表的可积层级（integrable hierarchy）有着惊人的联系。

% [Tutor Note: KdV 方程（Korteweg–de Vries）与 KP 方程（Kadomtsev–Petviashvili）都是可积的孤子方程；「可积层级」指由无穷多个彼此相容的方程构成的族。]

正如 2 维的引力理论是 4 维引力理论的玩具一样，2 维的规范理论是 4 维规范理论的玩具。

\takeaway{2 维既有玩具价值，又是复几何、代数几何与规范理论交汇之处。}

% [原文 p.120]
一般来说，引力理论比规范理论难得多。这一点在 2 维也成立，但引力理论比规范理论历史更久，自古以来就被研究，所以先稍微回忆一下 2 维经典纯引力理论。这里「纯（pure）」指不考虑引力以外的物质场、规范场等。

纯引力理论的运动方程由真空中的 Einstein 方程
\begin{equation}
  R_{\mu\nu} = \Lambda g_{\mu\nu}
  \tag{9.1.1}
\end{equation}
给出。这里 $R_{\mu\nu}$ 是 Ricci 曲率，$g_{\mu\nu}$ 是度量张量，$\Lambda$ 是宇宙常数（cosmological constant）。必要时把 Riemann 度量乘以常数，可以把 $\Lambda$ 规范化为 $0,1,-1$ 中的某一个。满足 Einstein 方程的 Riemann 流形称为 \textbf{Einstein 流形}（Einstein manifold）。在 2 维，Ricci 曲率可用标量曲率 $R$ 写成
\begin{equation}
  R_{\mu\nu} = R g_{\mu\nu}
  \tag{9.1.2}
\end{equation}
所以 (9.1.1) 成为
\begin{equation}
  R = \Lambda
  \tag{9.1.3}
\end{equation}
再者，曲率张量可用标量曲率写成
\begin{equation}
  R_{\mu\nu\rho\sigma} = R(g_{\mu\rho}g_{\nu\sigma} - g_{\mu\sigma}g_{\nu\rho})
  \tag{9.1.4}
\end{equation}
所以「是 Einstein 流形」与「是常曲率空间（space form）」等价。由 Killing--Hopf 定理，单连通的完备常曲率空间只限于球面、Euclid 空间或双曲空间，一般的完备常曲率空间都以其中之一为普遍覆叠。

% [Tutor Note: 2 维时曲率张量只剩一个独立分量，故 (9.1.4) 右端只有单一项 $R$；这是「Einstein $\Longleftrightarrow$ 常曲率」在 2 维成立的原因。]
% [Tutor Note: Killing–Hopf 定理（1898）分类单连通完备常曲率空间：$S^n$、$\R^n$、$\HH^n$ 三种模型空间。]

在 2 维，复几何、代数几何、微分几何以及群论之间的距离比其他维数近得多。实际上，下面 4 个概念在 2 维是一致的：
\begin{itemize}
  \item 紧复流形
  \item 代数函数域
  \item Riemann 流形的共形同值类
  \item 常曲率空间的自同构群的离散部分群中作用自由且余紧者
\end{itemize}
不过，这里总把流形连定向一起考虑。与其他维数相比，2 维的显著特长之一是共形结构与复结构一致。这来自两点：在 2 维，给定度量后，关于它旋转 $\pi/2$ 便确定近复结构；在 2 维，所有近复结构都可积。因此，共形映射与全纯映射一致，其自由度为无限维。这一点使共形不变量在 2 维格外强有力。

% [Tutor Note: 「近复结构可积」指存在相容的复坐标图册，即 Newlander–Nirenberg 条件 $\bar\partial^2=0$ 成立；2 维时该条件自动满足。]

% [原文 p.121]
在 2 维，对任意共形结构，存在唯一与之相容的常曲率 Riemann 度量。实际上，对任意 Riemann 面，考虑其普遍覆叠便得到单连通的 Riemann 面，而由 Poincaré--Koebe 单值化定理（uniformization theorem），单连通 Riemann 面必双全纯同构于
\begin{itemize}
  \item Riemann 球面 $\PP^1$
  \item 复平面 $\C$
  \item 上半平面 $\HH$
\end{itemize}
中的某一个。这些空间自然带有常曲率 Riemann 度量，所以把该双全纯同构像拉回，就得到与给定共形结构相容的常曲率 Riemann 度量。

与引力理论一样，在规范理论中，2 维时复几何、代数几何、微分几何与群论之间也有密切关系。特别地，下面 4 个概念在紧 Riemann 面上几乎一致：
\begin{itemize}
  \item 全纯向量丛
  \item 代数向量丛
  \item 常中心曲率联络
  \item 基本群（的中心扩张）的酉表示
\end{itemize}
这里说「几乎」，是因为忽略了稳定性问题。更准确的主张由 Narasimhan--Seshadri 的定理 9.10.1 给出，它构成本章的主定理。

% [Tutor Note: 常中心曲率联络在文献中多称 Hermite–Einstein 联络；「中心」指曲率取值为李代数中心里的标量矩阵。]

全纯向量丛是复几何的概念，代数向量丛是代数几何的概念。这两个概念的差别在于：用粘合构造向量丛时，变换函数取复解析意义下的全纯函数（holomorphic function），还是取代数几何意义下的正则函数（regular function）。前者指局部由收敛幂级数定义的函数，后者指局部由多项式定义的函数，复解析意义下的全纯函数概念比代数几何意义下的正则函数概念宽得多。作为其结果，任意代数向量丛自动也是全纯向量丛，但一般全纯向量丛未必是代数的。不过，由所谓 GAGA 原理，在射影流形上，全纯向量丛与代数向量丛成为等价的概念。这里 GAGA 这一名称来自 Serre 的论文 [122]。

% [Tutor Note: GAGA 是 Serre 论文标题「Géométrie Algébrique et Géométrie Analytique」的缩写，1956 年发表于 Ann. Inst. Fourier。]

常中心曲率联络是微分几何的概念，不用变换函数而用联络的语言定义。在理论物理中，联络称为规范场（gauge field），其曲率称为场强（field strength）。另外，把曲率的 $L^2$ 范数的平方对应到联络的函数称为 Yang--Mills 泛函。刻画这个 Yang--Mills 泛函的驻点的方程是 Yang--Mills 方程，它成为关于联络的 2 阶偏微分方程。这个 Yang--Mills 方程的解的模空间（充当规范等价类的空间）的研究，是规范理论的中心课题。在 2 维，正如 Einstein

% [原文 p.122]
方程与常曲率方程等价那样，Yang--Mills 方程也与常曲率方程
\begin{equation}
  *F = -2\pi\sqrt{-1}\Lambda
  \tag{9.1.5}
\end{equation}
等价。这里 $*$ 是 Hodge 星算子，$\Lambda$ 是酉群的李代数的元素。在 2 维，Yang--Mills 联络与常曲率联络一致，以及常曲率联络与基本群的中心扩张的酉表示对应，这些都很容易看出。这些事实的证明例如见 Atiyah--Bott [11]。这是关于 2 维规范理论的经典文献，把曲率理解为动量映射、把模空间看作无限维流形由无限维 Lie 群作出的辛商、对动量映射应用 Morse 理论等，其后成为标准的许多想法都在那里被引入。这篇论文篇幅很大，但写得极好，简直像教科书。若想介绍这篇论文的内容，也只能成为其劣化的复制，所以这里不作尝试。

% [Tutor Note: Atiyah–Bott（1982, Phil. Trans. Roy. Soc. A 308, 523–615）给出 9.5 节无限维辛商观点，是本章的理论源头。]

全纯向量丛与基本群的酉表示之间的关系，是单值化在规范理论中的对应物，由 Narasimhan--Seshadri [108] 在 Atiyah--Bott 之前就已知道，而 Atiyah--Bott 把 Narasimhan--Seshadri 的定理重新解释为无限维辛商。基于这种新解释给出 Narasimhan--Seshadri 定理的新证明这一问题，当时还是研究生的 Donaldson 从 Hitchin 那里得到。本章的目标是解说作为其结果得到的 Donaldson 对 Narasimhan--Seshadri 定理的微分几何别证（[35]）。

\section{向量丛的联络}

设 $E$ 为流形 $X$ 上的向量丛。线性映射
\begin{equation}
  d_A\colon \Gamma(E) \to \Omega^1(E)
  \tag{9.2.1}
\end{equation}
对任意 $f\in C^\infty(X)$、$s\in\Gamma(E)$ 满足 Leibniz 律
\begin{equation}
  d_A(fs) = f d_A(s) + s\otimes \rd f
  \tag{9.2.2}
\end{equation}
者称为 $E$ 的\textbf{联络}（connection）。$E$ 的联络所成的空间记作 $\cA(E)$。对任意 $A,A'\in\cA(E)$ 与 $t\in\R$，$(1-t)d_A + t d_{A'}$ 仍是 $E$ 的联络；又 $d_A - d_{A'}$ 给出 $\Omega^1(\End E)$ 的元素，因此可知 $\cA(E)$ 是以 $\Omega^1(\End E)$ 为模型的仿射空间。联络 $d_A\in\cA(E)$ 满足 Leibniz 律
\begin{equation}
  d_A(s\otimes \omega) = d_A(s)\wedge \omega + (-1)^{\deg s} s\otimes \rd\omega
  \tag{9.2.3}
\end{equation}

% [原文 p.123]
从而自然引起映射
\begin{equation}
  d_A\colon \Omega^\bullet(E) \to \Omega^{\bullet+1}(E)
  \tag{9.2.4}
\end{equation}
但这里对 $s\in\Omega^p(E)$ 设 $\deg s = p$。
\begin{align}
  d_A^2(fs) &= d_A(s\otimes \rd f + f d_A(s))
  \tag{9.2.5}\\
  &= d_A(s)\wedge \rd f + \rd f\wedge d_A(s) + f d_A^2(s)
  \tag{9.2.6}\\
  &= f d_A^2(s)
  \tag{9.2.7}
\end{align}
所以复合
\begin{equation}
  d_A^2 := d_A\circ d_A\colon \Omega^\bullet(E) \to \Omega^{\bullet+2}(E)
  \tag{9.2.8}
\end{equation}
不含微分。因此，存在某个 $\End E$ 值的 2 次微分形式
\begin{equation}
  F_A \in \Omega^2(\End E)
  \tag{9.2.9}
\end{equation}
使得 $d_A^2$ 成为由 $F_A$ 作乘法的算子；
\begin{equation}
  d_A^2(s) = F_A\wedge s.
  \tag{9.2.10}
\end{equation}
这个 $F_A$ 称为 $d_A$ 的\textbf{曲率}（curvature）。用局部标架写
\begin{equation}
  d_A = d + A
  \tag{9.2.11}
\end{equation}
则
\begin{align}
  F_A(s) &= (d+A)\wedge(d+A)s
  \tag{9.2.12}\\
  &= d^2 s + d(As) + A\wedge \rd s + A\wedge As
  \tag{9.2.13}\\
  &= 0 + ((dA)\wedge s - A\wedge \rd s) + A\wedge \rd s + A\wedge As
  \tag{9.2.14}\\
  &= (dA + A\wedge A)s
  \tag{9.2.15}
\end{align}
故
\begin{equation}
  F_A = dA + A\wedge A
  \tag{9.2.16}
\end{equation}

% [Tutor Note: (9.2.14) 的减号来自分次 Leibniz 律 $d(A\wedge s)=dA\wedge s+(-1)^{|A|}A\wedge \rd s$，此处 $|A|=1$。]

给定 $A\in\cA(E)$ 与 $a\in\Omega^1(\End E)$ 时，对任意 $s\in\Gamma(E)$，
\begin{align}
  F_{A+a}(s) &= (d_A+a)^2 s
  \tag{9.2.17}\\
  &= d_A^2 s + d_A(as) + a d_A s + a\wedge as
  \tag{9.2.18}\\
  &= F_A s + (d_A a)s - a\wedge d_A s + a\wedge d_A s + a\wedge as
  \tag{9.2.19}\\
  &= (F_A + d_A a + a\wedge a)s
  \tag{9.2.20}
\end{align}
故
\begin{equation}
  F_{A+a} = F_A + d_A a + a\wedge a
  \tag{9.2.21}
\end{equation}

\takeaway{联络是满足 Leibniz 律的协变导数，其曲率 $F_A=\rd A+A\wedge A$ 度量 $d_A^2\ne 0$ 的程度。}

\section{Yang--Mills 联络}

% [原文 p.124]
对 Riemann 流形 $X$ 上的主 $U(r)$ 丛，\textbf{Yang--Mills 泛函}（Yang--Mills functional）由
\begin{equation}
  \mathcal{YM}\colon \cA(P) \to \R,\qquad
  A \mapsto \|F_A\|^2_{L^2} = -\int_X \tr(F_A\wedge *F_A)
  \tag{9.3.1}
\end{equation}
定义。Yang--Mills 泛函的驻点称为 \textbf{Yang--Mills 联络}（Yang--Mills connection）。

\begin{theorem*}{9.3.1}
$A\in\cA(P)$ 为 Yang--Mills 联络的充要条件是
\begin{equation}
  d_A^* F_A = 0
  \tag{9.3.2}
\end{equation}
\end{theorem*}

\begin{proof}
对 $A\mapsto A+\delta A$ 有 $F_{A+\delta A} = F_A + d_A(\delta A)$，故
\begin{align}
  \delta\int_X \tr(F_A\wedge *F_A)
  \tag{9.3.3}\\
  &= \int_X \tr((\delta F_A)\wedge *F_A) + \int_X \tr(F_A\wedge *(\delta F_A))
  \tag{9.3.4}\\
  &= \int_X \tr(d_A(\delta A)\wedge *F_A) + \int_X \tr(F_A\wedge *(d_A\delta A))
  \tag{9.3.5}\\
  &= 2\int_X \tr(d_A(\delta A)\wedge *F_A)
  \tag{9.3.6}\\
  &= -2\int_X \tr(\delta A\wedge d_A(*F_A))
  \tag{9.3.7}
\end{align}
因此，注意 $d_A^* = *d_A *$ 便得 (9.3.2)。
\end{proof}

% [Tutor Note: (9.3.6) 用了 $X$ 上的内积 $\langle\alpha,\beta\rangle=\int\tr(\alpha\wedge*\beta)$ 的对称性与 $\tr$ 的循环性；(9.3.7) 用了一次分部积分并丢弃边界项。]

\takeaway{Yang--Mills 联络即 Yang--Mills 泛函的驻点，等价于 Yang--Mills 方程 $d_A^*F_A=0$。}

\section{二维流形上的拓扑向量丛}

设 $X$ 为 2 维流形，$E$ 为 $X$ 上的复向量丛。把 $E$ 的第 1 Chern 类在 $X$ 上积分所得的值称为 $E$ 的\textbf{次数}（degree）。由 Chern--Weil 理论，$E$ 具有联络 $A$ 时，可取 $A$ 的曲率的迹作为第 1 Chern 类在 de Rham 上同调群中的代表元；
\begin{equation}
  c_1(E) = \left[\frac{\sqrt{-1}}{2\pi}\tr F_A\right] \in H^2(X;\Z).
  \tag{9.4.1}
\end{equation}
2 维流形上的复向量丛的拓扑同型类如下分类：

% [原文 p.125]
\begin{proposition*}{9.4.1}
2 维流形上的复向量丛的拓扑同型类仅由秩与次数决定。
\end{proposition*}

\begin{proof}
位相空间 $X$ 上秩 $r$ 的复向量丛的拓扑同型类，由 $X$ 到 $BU(r)$ 的映射的同伦类分类，其中 $BU(r)$ 是 $U(r)$ 的分类空间（classifying space）。由胞体近似定理（cellular approximation theorem），CW 复体间的连续映射与胞体映射（即对任意 $i\in\N$ 把 $X$ 的 $i$ 骨架（$i$-skeleton）移到 $Y$ 的 $i$ 骨架的连续映射）同伦。取分类空间的模型
\[
  BU(r) = \varinjlim\bigl(\mathrm{Gr}(r,r)\hookrightarrow \mathrm{Gr}(r,r+1)\hookrightarrow \mathrm{Gr}(r,r+2)\hookrightarrow\cdots\bigr)
\]
则其 2 骨架是 $\PP^1\subset \mathrm{Gr}(r,r+1)\subset BU(r)$。从 2 维流形到 $\PP^1$ 的映射的同伦类仅由次数决定，故命题 9.4.1 成立。

另外，也可用如下直接进行障碍论式考察来证明：在 $X$ 的 0 骨架 $X^{(0)}$ 上取 $E$ 的平凡化，它显然可延拓为 $X$ 的 1 骨架 $X^{(1)}$ 上的平凡化。把它延拓到 $X$ 的 2 骨架 $X^{(2)}$ 上的平凡化的障碍就是 $E$ 的第 1 Chern 类。给定两个复向量丛 $E_1$ 与 $E_2$ 时，把它们在 $X^{(1)}$ 上平凡化，便可在 $X^{(1)}$ 上作出同构映射 $E_1|_{X^{(1)}}\xrightarrow{\sim} E_2|_{X^{(1)}}$。这个同构能延拓为 $X^{(2)}$ 上的同构的充要条件是 $E_1$ 与 $E_2$ 的第 1 Chern 类相等。之后注意 $X$ 为 2 维时 $X^{(2)}=X$，便得命题 9.4.1。
\end{proof}

% [Tutor Note: $\mathrm{Gr}(r,N)$ 为 $\C^N$ 的 $r$ 维复子空间构成的 Grassmann 流形；$\PP^1=\mathrm{Gr}(1,2)$ 给出 $BU(r)$ 的 2 骨架，故 2 维底流形上只需一个整值不变量。]

\takeaway{2 维流形上的复向量丛在同型意义下由秩与第 1 Chern 数（次数）完全决定。}

\section{作为动量映射的曲率}

设 $X$ 为 2 维流形，$E$ 为 $X$ 上的复向量丛。选取 $E$ 的一个 Hermite 度量 $h$ 并固定。选取 $E$ 的 Hermite 度量，等价于把伴随于 $E$ 的主 $GL_r(\C)$ 丛约化为主 $U(r)$ 丛。由 Cartan 分解或 Iwasawa 分解，$U(r)$ 与 $GL_r(\C)$ 同伦等价，所以这种约化必然存在。另外，与 Riemann 度量的情形一样，也可用 1 的分解具体作出 Hermite 度量。令 $G=U(r)$，则其复化 $G_\C$ 由 $GL(r,\C)$ 给出。

$E$ 的联络 $A$ 对任意 $s,t\in\Gamma(E)$ 满足
\begin{equation}
  \rd h(s,t) = h(d_A s, t) + h(s, d_A t)
  \tag{9.5.1}
\end{equation}
时，$A$ 称为\textbf{酉联络}（unitary connection）。$(E,h)$ 的酉联络全体所成的空间记作 $\cA(E,h)$。取 $(E,h)$ 所伴随的主 $G$ 丛为 $P$，则 $\cA=\cA(E,h)$ 自然与 $P$ 的联络空间 $\cA(P)$ 视为同一，

% [原文 p.126]
成为以 $\Omega^1(P\times_G\mathfrak{g})$ 为模型的仿射空间。

$(E,h)$ 作为 Hermite 向量丛的自同构群
\begin{equation}
  \cG(E,h) := \Aut(E,h)
  \tag{9.5.2}
\end{equation}
称为\textbf{规范变换群}（gauge transformation group）。规范变换群的李代数是 $\Lie\cG = \Gamma(P\times_G\mathfrak{g})$。由内积
\begin{equation}
  \Gamma(P\times_G\mathfrak{g})\times\Gamma(P\times_G\mathfrak{g})\to\R,\quad
  (f,g)\mapsto \int_X \tr(fg)\, \rd\vol_X
  \tag{9.5.3}
\end{equation}
可把 $(\Lie\cG)^*$ 与 $\Gamma(P\times_G\mathfrak{g})$ 视为同一。

对联络空间 $\cA$ 上的 2 次微分形式 $\omega$，取 $A\in\cA$ 与 $\alpha,\beta\in T_A\cA\cong\Omega^1(P\times_G\mathfrak{g})$，
\begin{equation}
  \omega_A(\alpha,\beta) = \int_X \tr(\alpha\wedge\beta)
  \tag{9.5.4}
\end{equation}
定义之，则由 $\mathfrak{g}$ 的内积 $\tr\colon \mathfrak{g}\otimes\mathfrak{g}\to\R$ 与 $X$ 上的积分 $\int_X\colon \Omega^1(X)\ni\eta,\tau\mapsto \int_X \eta\wedge\tau\in\R$ 的非退化性，$\omega$ 非退化。而且 (9.5.4) 关于 $A$ 是常数，所以作为无限维仿射空间 $\cA$ 上的微分形式，$\rd\omega=0$ 成立。

规范变换 $g\in\cG$ 对 $\cA$ 的作用，对 $A\in\cA$ 与 $\alpha\in\Omega^1(P\times_G\mathfrak{g})$，
\begin{align}
  d_{A+\alpha} = d_A+\alpha \mapsto d_{g(A+\alpha)} &= g\circ(d_A+\alpha)\circ g^{-1}
  \tag{9.5.5}\\
  &= g\circ d_A\circ g^{-1} + g\alpha g^{-1}
  \tag{9.5.6}\\
  &= d_{g(A)} + g\alpha g^{-1}
  \tag{9.5.7}
\end{align}
给出，所以 $\cG$ 对联络空间的作用 $g_*\colon T_A\cA\to T_{g(A)}\cA$ 是
\begin{equation}
  g_*\alpha = g\alpha g^{-1}
  \tag{9.5.8}
\end{equation}
并且由迹的循环不变性显然保持 $\omega$；
\begin{align}
  \omega_{g(A)}(g_*\alpha,g_*\beta) &= \int_X \tr(g\alpha g^{-1}\wedge g\beta g^{-1})
  \tag{9.5.9}\\
  &= \int_X \tr(\alpha\wedge\beta)
  \tag{9.5.10}\\
  &= \omega_A(\alpha,\beta).
  \tag{9.5.11}
\end{align}

另外，对 $h\in\Gamma(P\times_G\mathfrak{g})\cong\Lie(\cG)$，由 $e^{ht}\in\cG$ 引起的 $A\in\cA$ 的规范变换关于 $t$ 的 1 次项是
\begin{align}
  (1+th)(d_A) &= (1+th)\circ d_A\circ(1+th)^{-1}
  \tag{9.5.12}\\
  &= (1+th)\circ d_A\circ(1-th)
  \tag{9.5.13}\\
  &= (1+th)(1-th)d_A + (1+th)\circ d_A(-th)
  \tag{9.5.14}\\
  &= d_A - t\, d_A h
  \tag{9.5.15}
\end{align}

% [原文 p.127]
所以，$h$ 生成的 $\cA$ 上的向量场 $\underline{h}$ 在 $A\in\cA$ 处的值 $\underline{h}_A$ 由
\begin{equation}
  \underline{h}_A = -d_A h \in \Omega^1(P\times_G\mathfrak{g})
  \tag{9.5.16}
\end{equation}
给出。因此，对任意 $\alpha\in T_A\cA$，
\begin{align}
  \omega_A(\underline{h}_A,\alpha) &= -\int_X \tr((d_A h)\wedge\alpha)
  \tag{9.5.17}\\
  &= \int_X \tr(h\, d_A\alpha)
  \tag{9.5.18}\\
  &= \int_X \tr(hF_A)
  \tag{9.5.19}\\
  &= \int_X \tr(h*F_A)\, \rd\vol_X
  \tag{9.5.20}
\end{align}
这里用了：对 $s\in\Gamma(P\times_G\mathfrak{g})$ 与 $\eta\in\Omega^1(P\times_G\mathfrak{g})$，
\begin{align}
  0 &= \int_{\partial X}\tr(s\alpha)
  \tag{9.5.21}\\
  &= \int_X \rd\tr(s\alpha)
  \tag{9.5.22}\\
  &= \int_X \tr(d_A s\wedge\alpha) + \int_X \tr(s\, d_A\alpha)
  \tag{9.5.23}
\end{align}

% [Tutor Note: 原书此句引入 $\eta\in\Omega^1(P\times_G\mathfrak g)$，但下面 (9.5.21)–(9.5.23) 中只出现 $\alpha$；疑为笔误，照录。]

因此，映射
\begin{equation}
  *F\colon \cA \to (\Lie\cG)^*,\quad A\mapsto *F_A
  \tag{9.5.24}
\end{equation}
是 $\cG$ 对 $\cA$ 的作用的动量映射。酉群 $G$ 的李代数 $\mathfrak{g}$ 的中心由实标量矩阵构成。对实数 $\mu$，以 $\mu\cdot \id_E\in\Gamma(P\times_G\mathfrak{g})$ 为元的 1 点集合是 $\Lie(\cG)^*$ 的余伴随轨道。把满足微分方程
\begin{equation}
  *F_A = -2\pi\sqrt{-1}\mu\,\id_E
  \tag{9.5.25}
\end{equation}
的联络称为\textbf{常中心曲率联络}（constant central curvature connection），把常中心曲率联络所成的空间用规范变换群 $\cG$ 的作用切出的商空间
\begin{equation}
  \cM(r,d) := \bigl\{A\in\cA \bigm| *F_A = -2\pi\sqrt{-1}\mu\,\id_E\bigr\}\big/\cG
  \tag{9.5.26}
\end{equation}
称为常中心曲率联络的模空间。但这里
\begin{equation}
  d = r\mu
  \tag{9.5.27}
\end{equation}
另外，这是无限维辛仿射空间 $\cA$ 由无限维 Lie 群 $\cG$ 作出的辛商。$\cM(r,d)$ 一般有奇点，但与有限维辛商的情形一样，其非奇异部分自然带有辛结构。规范变换
\begin{equation}
  \cG\to\cA,\quad g\mapsto g(A)
  \tag{9.5.28}
\end{equation}

% [原文 p.128]
的微分 $\Lie(\cG)\to T_A\cA$ 由 (9.5.16) 给出为
\begin{equation}
  d_A\colon \Gamma(P\times_G\mathfrak{g})\to\Omega^1(P\times_G\mathfrak{g}),\quad
  h\mapsto -d_A h
  \tag{9.5.29}
\end{equation}
而把曲率对应到联络的映射
\begin{equation}
  F\colon \cA\to\Omega^2(P\times_G\mathfrak{g}),\quad A\mapsto F_A
  \tag{9.5.30}
\end{equation}
的微分由 (9.2.21) 给出为
\begin{equation}
  d_A\colon \Omega^1(P\times_G\mathfrak{g})\to\Omega^2(P\times_G\mathfrak{g}),\quad
  \alpha\mapsto d_A\alpha
  \tag{9.5.31}
\end{equation}
因此，$\cM(r,d)$ 的非奇异点 $[A]$ 处的切空间由复体
\begin{equation}
  0 \to \Gamma(P\times_G\mathfrak{g}) \xrightarrow{d_A}
  \Omega^1(P\times_G\mathfrak{g}) \xrightarrow{d_A}
  \Omega^2(P\times_G\mathfrak{g}) \to 0
  \tag{9.5.32}
\end{equation}
的 1 次上同调 $H^1(P\times_G\mathfrak{g})$ 给出。(9.5.4) 在辛商 $\cM_A$ 上诱导的辛形式，由上同调的杯积
\begin{equation}
  \cup\colon H^1(P\times_G\mathfrak{g})\times H^1(P\times_G\mathfrak{g})
  \to H^2(P\times_G(\mathfrak{g}\otimes\mathfrak{g})),
  \tag{9.5.33}
\end{equation}
系数取迹
\begin{equation}
  \tr\colon H^2(P\times_G(\mathfrak{g}\otimes\mathfrak{g}))\to H^2(X;\R),
  \tag{9.5.34}
\end{equation}
以及 Poincaré 对偶
\begin{equation}
  \int_X\colon H^2(X;\R)\xrightarrow{\sim}\R
  \tag{9.5.35}
\end{equation}
的合成给出。若只定义 $\cM(r,d)$ 非奇异部分上的 2 次微分形式 $\omega$，则无需引入无限维空间，但在有限维框架内证明 $\omega$ 闭并不一定容易。

\takeaway{曲率 $*F_A$ 是规范群作用的动量映射，模空间即其无限维辛商。}

\section{全纯向量丛}

设 $X$ 为复流形，$E$ 为 $X$ 上的复向量丛。$E$ 的变换函数为全纯函数时，$E$ 称为\textbf{全纯向量丛}（holomorphic vector bundle）。换言之，全纯向量丛指带有可积复结构的复向量丛。在复流形的定义中假定复结构的可积性，但在复向量丛的定义中并不假定复结构的可积性，请注意这一点。下面的定理是 Newlander--Nirenberg 定理对向量丛的类似物：

\begin{theorem*}{9.6.1}
对复流形 $X$ 上的复向量丛 $E$，给 $E$ 一个可积的复结构，等价于给出微分算子
\begin{equation}
  \bar\partial_E\colon \Gamma(E)\to\Omega^{0,1}(E)
  \tag{9.6.1}
\end{equation}

% [原文 p.129]
使得
\begin{equation}
  \bar\partial_E^2\colon \Gamma(E)\to\Omega^{0,2}(E)
  \tag{9.6.2}
\end{equation}
为零者。
\end{theorem*}

\begin{proof}
给定 $X$ 上的全纯向量丛 $E$ 时，对 $X$ 的各点，取该点的邻域 $U$ 与 $U$ 上的局部全纯标架 $(e_i)_{i=1}^r$，微分算子
\begin{equation}
  \bar\partial_E\colon \Gamma(E|_U)\to\Omega^{0,1}(E|_U)
  \tag{9.6.3}
\end{equation}
用 $s_1,\dots,s_r\in C^\infty(U)$ 写成
\begin{equation}
  s = \sum_{i=1}^r e_i s_i
  \tag{9.6.4}
\end{equation}
的 $U$ 上的 $E$ 的截面，对之
\begin{equation}
  \bar\partial_E s := \sum_{i=1}^r e_i\,\bar\partial s_i
  \tag{9.6.5}
\end{equation}
定义。为说明这样逐邻域定义的 $\bar\partial_E$ 粘合起来给出大域上的微分算子，只须证明这个定义不依赖于局部全纯标架的取法。为此，取 $U$ 上另一个局部全纯标架 $(e_j')_{j=1}^r$，由于 $(e_i)_{i=1}^r$ 与 $(e_j')_{j=1}^r$ 都是局部全纯标架，存在全纯的变换函数
\begin{equation}
  g = (g_{ij})_{i,j=1}^r \in \Gamma(\End E),\qquad \bar\partial g = 0
  \tag{9.6.6}
\end{equation}
使得
\begin{equation}
  e_j = \sum_{i=1}^r e_i' g_{ij},\qquad j=1,\dots,r
  \tag{9.6.7}
\end{equation}
成立。把 $s$ 关于新标架的表示写成
\begin{equation}
  s = \sum_{j=1}^r e_j' s_j'
  \tag{9.6.8}
\end{equation}
时，成分之间的关系为
\begin{equation}
  \sum_{i=1}^r e_i' s_i' = \sum_{j=1}^r e_j s_j = \sum_{i,j=1}^r e_i' g_{ij} s_j
  \tag{9.6.9}
\end{equation}
故
\begin{equation}
  s_i' = \sum_{j=1}^r g_{ij} s_j
  \tag{9.6.10}
\end{equation}
这时，

% [原文 p.130]
\begin{align}
  \sum_{i=1}^r e_i' \bar\partial s_i' &= \sum_{i,j=1}^r e_i'\bar\partial(g_{ij}s_j)
  \tag{9.6.11}\\
  &= \sum_{i,j=1}^r e_i'(\bar\partial g_{ij})s_j + \sum_{i,j=1}^r e_i' g_{ij}(\bar\partial s_j)
  \tag{9.6.12}\\
  &= \sum_{i,j=1}^r e_i' g_{ij}\bar\partial s_j
  \tag{9.6.13}\\
  &= \sum_{j=1}^r e_j \bar\partial s_j
  \tag{9.6.14}
\end{align}
所以可知 $\bar\partial_E$ 不依赖于局部全纯标架的取法而被定义。这样定义的 $\bar\partial_E$ 满足 $\bar\partial_E^2=0$，这由 $X$ 的复结构的可积性
\begin{equation}
  \bar\partial^2 = 0
  \tag{9.6.15}
\end{equation}
直接推出。

反之，关于 $\bar\partial_E^2=0$ 的微分算子
\begin{equation}
  \bar\partial_E\colon \Gamma(E)\to\Omega^{0,1}(E)
  \tag{9.6.16}
\end{equation}
在 $E$ 上确定可积复结构一事，其证明此处省略。例如见 [37, Theorem 2.1.53]。
\end{proof}

复流形 $X$ 上的复向量丛 $E$ 的联络
\begin{equation}
  d_A\colon \Gamma(E)\to\Omega^1(E)
  \tag{9.6.17}
\end{equation}
给定时，由 $X$ 的复结构定出的分解
\begin{equation}
  \Omega^1(X)\otimes_\R\C \cong \Omega^{1,0}(X)\oplus\Omega^{0,1}(X)
  \tag{9.6.18}
\end{equation}
可把 $d_A$ 分解为
\begin{equation}
  \partial_A\colon \Gamma(E)\to\Omega^{1,0}(E)
  \tag{9.6.19}
\end{equation}
与
\begin{equation}
  \bar\partial_A\colon \Gamma(E)\to\Omega^{0,1}(E)
  \tag{9.6.20}
\end{equation}
之和
\begin{equation}
  d_A = \partial_A + \bar\partial_A
  \tag{9.6.21}
\end{equation}

复流形 $X$ 上的全纯向量丛 $E$ 上的联络 $A$ 满足 $\bar\partial_A = \bar\partial_E$ 时，$A$ 称为与 $E$ 的全纯结构\textbf{相容}（compatible with the holomorphic structure of $E$）。

% [Tutor Note: 「与全纯结构相容」按定义即 $(0,1)$ 分量等于该全纯结构自带的 $\bar\partial_E$；这是后面标准联络存在唯一性的前提。]

\takeaway{全纯结构等价于满足 $\bar\partial_E^2=0$ 的 Dolbeault 算子；联络相容即 $\bar\partial_A=\bar\partial_E$。}

\section{标准联络}

% [原文 p.131]
下面的命题 9.7.1 中给出的联络称为\textbf{标准联络}（canonical connection）。

\begin{proposition*}{9.7.1}
对复流形 $X$ 上的全纯 Hermite 向量丛 $(E,h)$，与 $E$ 的全纯结构相容的酉联络唯一存在。
\end{proposition*}

命题 9.7.1 的证明大致分为用局部酉标架者与用局部全纯标架者两种，二者分别着眼于 Hermite 结构与全纯结构。

\begin{proof}[命题 9.7.1 的用局部酉标架的证明]
首先说明唯一性。对与 $E$ 的复结构相容的微分算子 $\bar\partial_E\colon \Gamma(E)\to\Omega^{0,1}(E)$，取 $(E,h)$ 的局部酉标架 $\{e_i\}_{i=1}^r$，把对应的 $(0,1)$ 形式记为 $\alpha=(\alpha_{ij})_{i,j=1}^r$；
\begin{equation}
  \bar\partial_E e_j = \sum_{i=1}^r e_i \alpha_{ij}.
  \tag{9.7.1}
\end{equation}
把关于 $(E,h)$ 的局部酉标架 $\{e_i\}_{i=1}^r$ 的 $E$ 的联络 $d_A$ 相伴的联络 1 形式记为 $A$，则由分解
\begin{equation}
  \Omega^1(X)\otimes\C = \Omega^{1,0}(X)\oplus\Omega^{0,1}(X)
  \tag{9.7.2}
\end{equation}
$A$ 分解为
\begin{equation}
  A = A^{1,0} + A^{0,1}
  \tag{9.7.3}
\end{equation}
由复共轭，1 次微分形式的 $(1,0)$ 部分与 $(0,1)$ 部分互换，所以
\begin{equation}
  (A^*)^{1,0} = (A^{0,1})^*,\qquad (A^*)^{0,1} = (A^{1,0})^*
  \tag{9.7.4}
\end{equation}
成立。是酉联络这一条件是
\begin{equation}
  A + A^* = 0
  \tag{9.7.5}
\end{equation}
而 $A$ 与 $E$ 的全纯结构相容这一条件是
\begin{equation}
  A^{0,1} = \alpha
  \tag{9.7.6}
\end{equation}
由这些条件得
\begin{equation}
  A^{1,0} = -(A^{0,1})^* = -\alpha^*,\qquad A = \alpha - \alpha^*
  \tag{9.7.7}
\end{equation}
因此，与 $E$ 的全纯结构相容的酉联络若存在则唯一。

其次说明存在性。选取给出 $E$ 的局部平凡化的 $X$ 的开覆盖 $\{U_\lambda\}_{\lambda\in\Lambda}$ 并固定之。对任意 $\lambda\in\Lambda$，在 $E|_{U_\lambda}$ 上存在由 (9.7.7) 确定的、与全纯结构相容的酉联络 $A_\lambda$。对任意 $\lambda,\mu\in\Lambda$，$A_\lambda|_{U_\lambda\cap U_\mu}$ 与 $A_\mu|_{U_\lambda\cap U_\mu}$ 都给出与全纯结构相容的 $E|_{U_\lambda\cap U_\mu}$ 的酉联络。由上面证明的唯一性，二者一致，所以这些局部联络粘合起来，确定 $E$ 上的整体联络。
\end{proof}

% [Tutor Note: (9.7.7) 表明标准联络的 1 形式完全由 $\alpha=\bar\partial$ 的标架矩阵决定，故又称 Chern 联络。]

% [原文 p.132]
\begin{proof}[命题 9.7.1 的用局部全纯标架的证明]
为说明唯一性，假设存在与 $X$ 以及 $E$ 的复结构相容的 $(E,h)$ 上的酉联络 $d_A$。选取 $E$ 的全纯局部标架 $(e_i)_{i=1}^r$，把与它相伴的联络 1 形式记为 $\omega=(\omega_{ij})_{i,j=1}^r$；
\begin{equation}
  d_A e_j = \sum_{i=1}^r e_i \omega_{ij}.
  \tag{9.7.8}
\end{equation}
由与全纯结构的相容性以及 $(e_i)_{i=1}^r$ 是全纯局部标架，
\begin{equation}
  \bar\partial_A e_i = \bar\partial_E e_i = 0
  \tag{9.7.9}
\end{equation}
故知 $\omega$ 是 $(1,0)$ 形式。
\begin{equation}
  h_{ij} = h(e_i,e_j)
  \tag{9.7.10}
\end{equation}
而 $d_A$ 保持 $h$，故
\begin{align}
  dh_{ij} &= d\bigl(h(e_i,e_j)\bigr)
  \tag{9.7.11}\\
  &= h(d_A e_i, e_j) + h(e_i, d_A e_j)
  \tag{9.7.12}\\
  &= h\Bigl(\sum_{k=1}^r e_k\omega_{ki},\, e_j\Bigr)
     + h\Bigl(e_i,\, \sum_{k=1}^r e_k\omega_{kj}\Bigr)
  \tag{9.7.13}\\
  &= \sum_{k=1}^r \omega_{ki} h_{kj} + \sum_{k=1}^r h_{ik}\bar\omega_{kj}
  \tag{9.7.14}
\end{align}
由于 $\omega$ 是 $(1,0)$ 形式，把 (9.7.14) 分成 $(1,0)$ 成分与 $(0,1)$ 成分便得
\begin{align}
  \partial h_{ij} &= \sum_{k=1}^r \omega_{ki} h_{kj},
  \tag{9.7.15}\\
  \bar\partial h_{ij} &= \sum_{k=1}^r h_{ik}\bar\omega_{kj}
  \tag{9.7.16}
\end{align}
用矩阵记号表示，则
\begin{align}
  \partial h &= \omega^T h,
  \tag{9.7.17}\\
  \bar\partial h &= h\bar\omega
  \tag{9.7.18}
\end{align}
由 (9.7.17)，

% [原文 p.133]
\begin{equation}
  \omega = \bigl((\partial h)h^{-1}\bigr)^T
  \tag{9.7.19}
\end{equation}
所以可知，与 $E$ 的全纯结构相容的 $(E,h)$ 上的酉联络若存在则唯一。
另一方面，由 $h$ 的 Hermite 性，对 (9.7.17) 两边取 Hermite 共轭便得 (9.7.18)。换言之，由 (9.7.19)，与全纯结构相容的酉联络在局部存在。此处再用唯一性，便知它们粘合起来确定整体的联络。
\end{proof}

\takeaway{与全纯结构相容的酉联络存在且唯一，即标准联络，其 1 形式 $\omega=((\partial h)h^{-1})^T$。}

\section{复规范变换群}

把复向量丛 $E$ 作为复向量丛的自同构群
\begin{equation}
  \cG_\C(P) := \Aut E
  \tag{9.8.1}
\end{equation}
称为\textbf{复规范变换群}（complexified gauge transformation group）。它是规范变换群的复化。
把 $P$ 的酉联络全体所成之空间记为 $\cA(P)$。对于联络 $A\in\cA(P)$，其曲率 $F_A\in\Omega^2(P\times_G\mathfrak{g})$ 按分解
\begin{equation}
  \Omega^2(X)\otimes_\R\C
  = \Omega^{2,0}(X)\oplus\Omega^{1,1}(X)\oplus\Omega^{0,2}(X)
  \tag{9.8.2}
\end{equation}
而分解为
\begin{equation}
  F_A = F_A^{2,0} + F_A^{1,1} + F_A^{0,2}
  \tag{9.8.3}
\end{equation}
并且由 $A$ 的酉性得
\begin{equation}
  \bigl(F_A^{2,0}\bigr)^* = -F^{0,2},\qquad
  \bigl(F_A^{1,1}\bigr)^* = -F^{1,1}
  \tag{9.8.4}
\end{equation}
% [Tutor Note: (9.8.4) 右端的 $F^{0,2},F^{1,1}$ 原书未写下标 $A$，照录；实即 $F_A^{0,2},F_A^{1,1}$。]
若置
\begin{equation}
  \cA^{1,1}(P) := \bigl\{A\in\cA(P)\bigm| F_A^{0,2} = 0\bigr\}
  \tag{9.8.5}
\end{equation}
则由定理 9.6.1，对任意 $A\in\cA^{1,1}(P)$，$\bar\partial_A$ 在 $E$ 上确定一个全纯结构。$A,A'\in\cA^{1,1}(P)$ 所确定的全纯向量丛 $\cE_A$ 与 $\cE_{A'}$ 同构的充要条件是：存在 $E$ 作为复向量丛的自同构 $g\in\Aut(E)$，使得对任意 $s\in\Gamma(E)$ 有
\begin{equation}
  (\bar\partial_{A'}\circ g)(s) = (g\circ\bar\partial_A)(s)
  \tag{9.8.6}
\end{equation}
成立。这可以改写为
\begin{equation}
  \bar\partial_{A'} = g\circ\bar\partial_A\circ g^{-1}
  = \bar\partial_A - (\bar\partial_A g)g^{-1}
  \tag{9.8.7}
\end{equation}

% [原文 p.134]
由命题 9.7.1，给出 $P$ 的酉联络 $d_A$ 与给出其 $(0,1)$ 成分 $\bar\partial_A$ 是等价的。对于联络 $A\in\cA^{1,1}$ 与复规范变换 $g\in\cG_\C$，把新的联络 $g(A)$ 定义为这样的酉联络：其 $(0,1)$ 成分满足
\begin{equation}
  \bar\partial_{g(A)} = \bar\partial_A - (\bar\partial_A g)g^{-1}
  \tag{9.8.8}
\end{equation}
此时 $g(A)\in\cA^{1,1}$ 的 $(1,0)$ 成分 $\partial_{g(A)}$，由 (9.7.7)，对 $\tilde g=(g^*)^{-1}$ 由
\begin{align}
  \partial_{g(A)} &= \partial_A + \bigl((\bar\partial_A g)g^{-1}\bigr)^*
  \tag{9.8.9}\\
  &= \partial_A + (g^*)^{-1}(\partial_A g^*)
  \tag{9.8.10}\\
  &= \tilde g\circ\partial_A\circ\tilde g^{-1}
  \tag{9.8.11}
\end{align}
给出。由此，$\cG_\C$ 在 $\cA^{1,1}$ 上的作用被确定。全纯向量丛的同型类所成之集 $\cM$，与 $\cA^{1,1}$ 的 $\cG_\C$ 轨道之集一一对应；
\begin{equation}
  \cM \cong \cA^{1,1}/\cG_\C.
  \tag{9.8.12}
\end{equation}

给定全纯 Hermite 向量丛 $(\cE,h)$，把 $(\cE,h)$ 的标准联络在 $\cG_\C$ 作用下的轨道记为 $O(\cE)$。全纯向量丛 $\cE$ 上 Hermite 结构的同型类所成之集，与 $O(\cE)$ 的 $\cG$ 轨道之集 $O(\cE)/\cG$ 一一对应。

上面固定了 $E$ 上的 Hermite 结构 $h$ 而让全纯结构 $\bar\partial$ 变动，但由 (9.7.18) 有
\begin{equation}
  \bar\omega = h^{-1}\bar\partial h
  \mapsto h^{-1}\bigl(g\bar\partial g^{-1}\bigr)h
  = (g^{-1}h)^{-1}\,\bar\partial\,(g^{-1}h)
  \tag{9.8.13}
\end{equation}
所以也可以看作固定全纯结构 $\bar\partial$ 而按
\begin{equation}
  h \mapsto g^{-1}h
  \tag{9.8.14}
\end{equation}
变动 Hermite 结构 $h$。这两种看法（即固定 Hermite 结构 $h$ 而变动全纯结构 $\bar\partial$，还是固定全纯结构 $\bar\partial$ 而变动 Hermite 结构 $h$）完全等价。以下在陈述定理的主张时采用后者，在作证明时采用前者。

\takeaway{复规范变换群作用在 $\cA^{1,1}$ 上，全纯丛的同型类与轨道空间 $\cA^{1,1}/\cG_\C$ 一一对应。}

\section{梯度与稳定性}

Riemann 面上的向量丛，把次数与秩之比称为\textbf{梯度}（slope）；
\begin{equation}
  \mu(\cE) = \frac{\deg\cE}{\rank\cE}.
  \tag{9.9.1}
\end{equation}
全纯向量丛 $\cE$ 若对任意全纯子丛 $\cF$ 都有
% [原文 p.135]
\begin{equation}
  \mu(\cF) \le \mu(\cE)
  \tag{9.9.2}
\end{equation}
则称为\textbf{半稳定}（semistable）。更强地，若对任意真全纯子丛 $\cF$ 都有
\begin{equation}
  \mu(\cF) < \mu(\cE)
  \tag{9.9.3}
\end{equation}
则称 $\cE$ 是\textbf{稳定}（stable）的。又，非半稳定的全纯向量丛称为\textbf{不稳定}（unstable）。全纯向量丛 $\cE$ 称为\textbf{可分解}（decomposable），指存在适当的向量丛 $\cE_1,\cE_2$ 使得 $\cE\cong\cE_1\oplus\cE_2$。非可分解的全纯向量丛称为\textbf{直和不可分解}（indecomposable）。对酉联络也同样，与两个酉联络的直和同型时称为\textbf{可约}（reducible），不可约时称为\textbf{不可约}（irreducible）。

% [Tutor Note: 原书把向量丛的「直既約」gloss 作 indecomposable，把联络的「直既約」gloss 作 irreducible，同一词形两种含义；译文分别作「直和不可分解」与「不可约」。]

\takeaway{以梯度 $\mu=\deg\cE/\rank\cE$ 比较子丛与全丛，即可定义半稳定与稳定。}

\section{Narasimhan--Seshadri 定理}

把定中心曲率联络的方程 (9.5.25) 两边取迹，在 $X$ 上积分，便得
\begin{equation}
  \deg\cE = \mu\rank\cE
  \tag{9.10.1}
\end{equation}
即
\begin{equation}
  \mu = \mu(\cE)
  \tag{9.10.2}
\end{equation}
这里假设 $X$ 的体积已规一化为 $1$。特别地，对于次数为 $0$ 的向量丛，具有常中心曲率的联络与曲率为 $0$、即成为平坦联络（flat connection）是等价的。

\begin{theorem*}{9.10.1}
（Narasimhan--Seshadri [108]）Riemann 面 $X$ 上的直和不可分解全纯向量丛 $\cE$ 稳定的充要条件是：存在 $\cE$ 上的 Hermite 度量，使与它相伴的酉联络具有常中心曲率。又，对任意直和不可分解全纯向量丛 $\cE$，其上的常中心曲率 Hermite 度量若存在则唯一。
\end{theorem*}

以下利用本章余下的部分，介绍 Donaldson [35] 对 Narasimhan--Seshadri 定理的证明。

\takeaway{稳定 $\Leftrightarrow$ 存在常中心曲率的 Hermite 度量（Narasimhan--Seshadri），且该度量唯一。}

\section{Sobolev 空间}

本节复习 Sobolev 空间的若干基本事实。设 $X$ 为 $n$ 维紧 Riemann 流形。测度空间上的函数 $f$ 与 $g$，在除去零集后一致时定义为等价。对 $p\in[1,\infty)$，把使得 $|f|^p$ 可积的 $X$ 上可测函数的等价类所成之集记为 $L^p(X)$。

% [原文 p.136]
又，对 $p=\infty$，把\textbf{本质有界}（essentially bounded，即除去零集后有界）的可测函数的等价类所成之集记为 $L^\infty(X)$。函数 $f\colon X\to\R\cup\{\infty\}$ 的\textbf{本质上确界}（essential supremum）由
\begin{equation}
  \operatorname*{ess\,sup}_{x\in X} f
  := \inf\bigl\{a\in\R\cup\{\infty\}
  \bigm| \text{对几乎所有 } x\in X \text{ 有 } f(x)\le a\bigr\}
  \tag{9.11.1}
\end{equation}
定义。$L^p(X)$ 成为以
\begin{equation}
  \|f\|_{L^p(X)} :=
  \begin{cases}
    \Bigl(\displaystyle\int_X |f|^p \rd\vol_X\Bigr)^{\frac1p},
      & p\in[1,\infty),\\[2mm]
    \displaystyle\operatorname*{ess\,sup}_{x\in X} |f(x)|,
      & p=\infty
  \end{cases}
  \tag{9.11.2}
\end{equation}
为范数的 Banach 空间。$L^p$ 范数的三角不等式称为 \textbf{Minkowski 不等式}（Minkowski's inequality）。对 $L^p$ 空间，下面的 \textbf{Hölder 不等式}（Hölder's inequality）是基本的：

\takeaway{$L^p$ 是模去零集差异的 Banach 空间，其基本工具是 Hölder 与 Minkowski 不等式。}

\begin{theorem*}{9.11.1}
（Hölder 不等式）若 $p,q,r\in[1,\infty]$ 满足
\begin{equation}
  \frac1p + \frac1q = \frac1r
  \tag{9.11.3}
\end{equation}
则对任意 $f\in L^p(X)$ 与 $g\in L^q(X)$ 有
\begin{equation}
  \|fg\|_{L^r} \le \|f\|_{L^p}\|g\|_{L^q}
  \tag{9.11.4}
\end{equation}
成立。
\end{theorem*}

Minkowski 不等式可由 Hölder 不等式导出。由 $X$ 的紧性，对任意 $1\le p\le q\le\infty$ 有
\begin{equation}
  L^p(X) \supset L^q(X)
  \tag{9.11.5}
\end{equation}
成立。若 $p,p^*\in[1,\infty]$ 满足
\begin{equation}
  \frac1p + \frac1{p^*} = 1
  \tag{9.11.6}
\end{equation}
则称 $p$ 与 $p^*$ 为 \textbf{Hölder 共轭}（Hölder conjugate）。当 $p\in(1,\infty)$ 时，$L^{p^*}(X)$ 成为 $L^p(X)$ 作为拓扑线性空间的对偶空间。因此，对 $1<p<\infty$，$L^p(X)$ 是自反 Banach 空间。另一方面，$L^1(X)$ 的对偶空间是 $L^\infty(X)$，但 $L^\infty(X)$ 的对偶空间真包含 $L^1(X)$，所以 $L^1(X)$ 与 $L^\infty(X)$ 不是自反的。

对 Banach 空间 $V$ 的点列 $\{v_i\}_{i=1}^\infty$，若存在某个元 $v\in B$，使得对任意 $v^*\in V^*$，$\R$ 的点列 $\{v^*(v_i)\}_{i=1}^\infty$ 都收敛于 $v^*(v)$，则称 $\{v_i\}_{i=1}^\infty$ \textbf{弱收敛}（weakly converge）于 $v$。

% [Tutor Note: 上段「$v\in B$」原书如此，按上下文应为 $v\in V$（$B$ 疑为 $V$ 之误），照录。]

% [原文 p.137]
已知 Banach 空间的范数关于弱收敛是下半连续的；
\begin{equation}
  \|v\| \le \liminf_{i\to\infty}\|v_i\|.
  \tag{9.11.7}
\end{equation}
在 (9.11.7) 中等号不成立的典型例子，是平方可和点列所成的 Hilbert 空间
\begin{equation}
  \ell^2 := \Bigl\{(a_i)_{i=0}^\infty\in\C^\N
  \Bigm| \sum_{i=1}^\infty |a_i|^2 < \infty\Bigr\}
  \tag{9.11.8}
\end{equation}
中的正规正交基 $\{e_i\}_{i=1}^\infty$。事实上，此时 $\{e_i\}_{i=1}^\infty$ 当 $i\to\infty$ 时弱收敛于 $0$，故 (9.11.7) 的左边为 $0$，右边为 $1$。

$L^1(X)$ 的元 $f$，若存在某个 $g\in L^1(X,T^*X)$，使得对任意 $Y\in C^\infty(X,TX)$ 与任意 $\omega\in C^\infty(X,\wedge^nT^*X)$ 都有
\begin{equation}
  \int_X (\iota_Y g)\omega = -\int_X f(\cL_Y\omega)
  \tag{9.11.9}
\end{equation}
成立，则称 $f$ 为\textbf{弱可微}（weakly differentiable）。此时 $g$ 称为 $f$ 的\textbf{弱导数}（weak derivative）。对自然数 $k$ 与 $p\in[1,\infty]$，把 $k$ 阶弱可微、而且 $k$ 阶弱导数仍属于 $L^p(X)$ 的元全体所成的 $L^p(X)$ 的子空间称为 \textbf{Sobolev 空间}（Sobolev space），记作 $L^p_k(X)$。Sobolev 空间关于 \textbf{Sobolev 范数}（Sobolev norm）
\begin{equation}
  \|f\|_{L^p_k(X)} = \sum_{i=0}^k \|\nabla^i f\|_{L^p(X)}
  \tag{9.11.10}
\end{equation}
成为 Banach 空间。又，推广这一定义，对任意 $k\in\R$ 与 $p\in[1,\infty]$ 也可定义 $L^p_k(X)$。

\begin{theorem*}{9.11.2}
（Sobolev 嵌入定理，Sobolev embedding theorem）若 $k\ge l$ 且 $k-\dfrac np \ge l-\dfrac nq$，则 $L^p_k(X)\subset L^q_l(X)$。
\end{theorem*}

这里，对拓扑向量空间 $V$ 与 $W$，$V\subset W$ 指存在作为向量空间的嵌入，且该嵌入连续。因此，若 $k\ge l$ 且 $k-\dfrac np\ge l-\dfrac nq$，则对任意 $f\in L^p_k(X)$，$f$ 的 $L^p_k$ 范数被 $f$ 的 $L^q_l$ 范数从上控制。这称为 \textbf{Sobolev 不等式}（Sobolev inequality）。

% [Tutor Note: 原书称 $L^p_k$ 范数被 $L^q_l$ 范数从上控制，方向似与嵌入 $L^p_k\subset L^q_l$ 相反（通常为 $\|f\|_{L^q_l}\le C\|f\|_{L^p_k}$）；照录原书。]

Banach 空间之间的有界线性算子 $B\colon V\to W$，当 $V$ 的任意有界点列在 $B$ 下的像都具有收敛子列时称为\textbf{紧}（compact）的。当 $V$ 是 $W$ 的子空间且嵌入映射紧时，记作 $V\Subset W$。Sobolev 嵌入只要不是临界情形便是紧的：

\begin{theorem*}{9.11.3}
（Rellich--Kondrachov 定理）若 $k>l$ 且 $k-\dfrac np > l-\dfrac nq$，则 $L^p_k(X)\Subset L^q_l(X)$。
\end{theorem*}

% [原文 p.138]
又，下面的定理也称为 Sobolev 嵌入定理。

\begin{theorem*}{9.11.4}
若 $k-\dfrac np > l$，则 $L^p_k(X)\subset C^l(X)$。
\end{theorem*}

由 Sobolev 嵌入定理与 Hölder 不等式可得下面的事实：

\begin{theorem*}{9.11.5}
若 $k-\dfrac np>0$、$k\ge l$、$k-\dfrac np\ge l-\dfrac nq$ 且 $l\in\Z^{\ge0}$，则乘法映射
\begin{equation}
  L^p_k(X)\otimes L^q_l(X)\to L^q_l(X),\qquad f\otimes g\mapsto fg
  \tag{9.11.11}
\end{equation}
连续。
\end{theorem*}

\begin{proof}
由 Sobolev 范数的定义 (9.11.10)，
\begin{align}
  \|fg\|_{L^q_l} &= \sum_{i=0}^l \bigl\|\nabla^i(fg)\bigr\|_{L^q}
  \tag{9.11.12}\\
  &= \sum_{i=0}^l \Bigl\|\sum_{r=0}^i \binom{i}{r}(\nabla^r f)(\nabla^{i-r}g)\Bigr\|_{L^q}
  \tag{9.11.13}\\
  &\le \sum_{i=0}^l\sum_{r=0}^i \binom{i}{r}
    \bigl\|(\nabla^r f)(\nabla^{i-r}g)\bigr\|_{L^q}
  \tag{9.11.14}
\end{align}
故只需对任意满足 $a+b\le l$ 的自然数 $a,b$ 证明不等式
\begin{equation}
  \bigl\|(\nabla^a f)(\nabla^b g)\bigr\|_{L^q}
  \le C\|f\|_{L^p_k}\|g\|_{L^q_l}
  \tag{9.11.15}
\end{equation}
即可。这里 $C$ 是仅依赖于 $X,p,q,k,l,a,b$ 的常数。由 $f\in L^p_k(X)$ 得 $\nabla^a f\in L^p_{k-a}(X)$，而由 Sobolev 嵌入定理
\begin{equation}
  L^p_{k-a}(X)\subset
  \begin{cases}
    C^0(X), & (k-a)-\dfrac np > 0,\\[1mm]
    L^\infty(X), & (k-a)-\dfrac np = 0,\\[1mm]
    L^{p'}(X), & (k-a)-\dfrac np < 0
  \end{cases}
  \tag{9.11.16}
\end{equation}
成立。这里当 $(k-a)-\dfrac np<0$ 时，$p'\in[1,\infty)$ 由
\begin{equation}
  (k-a)-\frac np = -\frac n{p'}
  \tag{9.11.17}
\end{equation}
定义。由 $a\le l\le k$，有 $p<p'<\infty$。同样地，
\begin{equation}
  \nabla^b g\in L^q_{l-b}(X)\subset
  \begin{cases}
    C^0(X), & (l-b)-\dfrac nq > 0,\\[1mm]
    L^\infty(X), & (l-b)-\dfrac nq = 0,\\[1mm]
    L^{q'}, & (l-b)-\dfrac nq < 0
  \end{cases}
  \tag{9.11.18}
\end{equation}
成立。不过当 $(l-b)-\dfrac nq<0$ 时置
\begin{equation}
  (l-b)-\frac nq = -\frac n{q'}
  \tag{9.11.19}
\end{equation}
下面按 $(k-a)-\dfrac np$ 与 $(l-b)-\dfrac nq$ 的正负分类讨论：

% [原文 p.139]
\begin{enumerate}[label=\arabic*., leftmargin=2.2em]
\item 当 $(k-a)-\dfrac np\ge0$ 且 $(l-b)-\dfrac nq\ge0$ 时，
\begin{align}
  L^p_{k-a}(X)\otimes L^q_{l-b}(X)
  &\hookrightarrow L^\infty(X)\otimes L^\infty(X)
  \tag{9.11.20}\\
  &\hookrightarrow L^{2q}(X)\otimes L^{2q}(X)
  \tag{9.11.21}\\
  &\to L^q(X)
  \tag{9.11.22}
\end{align}
\item 当 $(k-a)-\dfrac np\ge0$ 且 $(l-b)-\dfrac nq<0$ 时，再作如下分类讨论：
  \begin{enumerate}[label=(\roman*), leftmargin=2.4em]
  \item 当 $b=l$ 时 $a=0$。由假设 $k-\dfrac np>0$，故 $L^p_k(X)\subset C^0(X)$，于是
\begin{align}
  L^p_{k-a}(X)\otimes L^q_{l-b}(X) &= L^p_k(X)\otimes L^q(X)
  \tag{9.11.23}\\
  &\hookrightarrow C^0(X)\otimes L^q(X)
  \tag{9.11.24}\\
  &\to L^q(X)
  \tag{9.11.25}
\end{align}
  \item 当 $b<l$ 时，令
\begin{equation}
  r = \frac{n}{l-b}
  \tag{9.11.26}
\end{equation}
则 $(l-b)-\dfrac nq<0$ 给出 $r>q$，并且
\begin{equation}
  \frac1r + \frac1{q'}
  = \frac{l-b}{n} + \Bigl(-\frac{l-b}{n} + \frac1q\Bigr)
  = \frac1q
  \tag{9.11.27}
\end{equation}
成立。由 $(k-a)-\dfrac np\ge0$ 得 $L^p_{k-a}(X)\subset L^\infty(X)$，故
\begin{align}
  L^p_{k-a}(X)\otimes L^q_{l-b}(X)
  &\hookrightarrow L^\infty(X)\otimes L^{q'}(X)
  \tag{9.11.28}\\
  &\hookrightarrow L^r(X)\otimes L^{q'}(X)
  \tag{9.11.29}\\
  &\to L^q(X)
  \tag{9.11.30}
\end{align}
  \end{enumerate}
\item 当 $(k-a)-\dfrac np<0$ 且 $(l-b)-\dfrac nq\ge0$ 时，
\begin{equation}
  (k-a)-\frac np \ge (k-l)-\frac np \ge -\frac nq
  \tag{9.11.31}
\end{equation}
故 $L^p_{k-a}\hookrightarrow L^q$。
  \begin{enumerate}[label=(\roman*), leftmargin=2.4em]
  \item 当 (9.11.31) 中取等号时，$a=l$ 且 $k-\dfrac np = l-\dfrac nq$。由 $a=l$ 得 $b=0$，而
\begin{equation}
  (l-b)-\frac nq = l-\frac nq = k-\frac np
  \tag{9.11.32}
\end{equation}

% [原文 p.140]
由假设 $k-\dfrac np>0$，故 $l-\dfrac nq>0$。于是 $L^q_l(X)\subset C^0(X)$，且
\begin{align}
  L^p_{k-a}(X)\otimes L^q_{l-b}(X)
  &\hookrightarrow L^q(X)\otimes C^0(C)
  \tag{9.11.33}\\
  &\to L^q(X)
  \tag{9.11.34}
\end{align}
成立。

% [Tutor Note: (9.11.33) 右端原书作 $C^0(C)$，按上下文应为 $C^0(X)$（$C$ 疑为 $X$ 之误），照录。]
  \item 当 (9.11.31) 中不取等号时，对充分小的正数 $\epsilon$ 有
\begin{equation}
  L^p_{k-a}\hookrightarrow L^{q+\epsilon}
  \tag{9.11.35}
\end{equation}
其中由
\begin{equation}
  \frac1{q+\epsilon} + \frac1r = \frac1q
  \tag{9.11.36}
\end{equation}
定义 $r\in(1,\infty)$，则
\begin{align}
  L^p_{k-a}(X)\otimes L^q_{l-b}(X)
  &\hookrightarrow L^{q+\epsilon}(X)\otimes L^\infty(X)
  \tag{9.11.37}\\
  &\hookrightarrow L^{q+\epsilon}(X)\otimes L^r(X)
  \tag{9.11.38}\\
  &\to L^q
  \tag{9.11.39}
\end{align}
成立。
  \end{enumerate}
\item 当 $(k-a)-\dfrac np<0$ 且 $(l-b)-\dfrac nq<0$ 时，由
\begin{equation}
  \frac1r = \frac1{p'} + \frac1{q'}
  \tag{9.11.40}
\end{equation}
定义 $r\in(0,\infty)$，则
\begin{align}
  \frac1r &= -\frac1n\Bigl[\Bigl\{(k-a)-\frac np\Bigr\}
    + \Bigl\{(l-b)-\frac nq\Bigr\}\Bigr]
  \tag{9.11.41}\\
  &= -\frac1n\Bigl[k+(l-a-b)-\frac np-\frac nq\Bigr]
  \tag{9.11.42}\\
  &\le -\frac1n\Bigl[k-\frac np-\frac nq\Bigr]
  \tag{9.11.43}\\
  &\le -\frac1n\Bigl[-\frac nq\Bigr]
  \tag{9.11.44}\\
  &= \frac1q
  \tag{9.11.45}
\end{align}
故
\begin{align}
  L^p_{k-a}(X)\otimes L^q_{l-b}(X)
  &\hookrightarrow L^{p'}(X)\otimes L^{q'}(X)
  \tag{9.11.46}\\
  &\to L^r(X)
  \tag{9.11.47}\\
  &\hookrightarrow L^q
  \tag{9.11.48}
\end{align}
成立。
\end{enumerate}
\end{proof}

